% Part: incompleteness
% Chapter: representability-in-q
% Section: introduction

\documentclass[../../../include/open-logic-section]{subfiles}

\begin{document}

\olfileid[hi]{inc}{req}{int}
\olsection{परिचय}

अपूर्णता-प्रमेय उन सिद्धांतों पर लागू होते हैं जिनमें संगणनीय फलनों के
मूलभूत तथ्यों को व्यक्त और सिद्ध किया जा सकता है। हम ऐसे एक अत्यंत न्यूनतम
सिद्धांत का वर्णन करेंगे, जिसे ``$\Th{Q}$'' (अथवा कभी-कभी रफ़ाएल रॉबिन्सन
के नाम पर ``रॉबिन्सन का $Q$'') कहते हैं। हम बताएँगे कि किसी फलन का
$\Th{Q}$ में \emph{निरूपणीय} होने का क्या अर्थ है और फिर यह सिद्ध करेंगे:
\begin{quote}
  कोई फलन $\Th{Q}$ में तभी और केवल तभी निरूपणीय है जब वह संगणनीय है।
\end{quote}
एक बात तो यह है कि इससे हमें संगणनीयता का एक और प्रतिरूप मिलता है। किंतु
हम इसका उपयोग यह दिखाने के लिए भी करेंगे कि समुच्चय
$\Setabs{!A}{\Th{Q} \Proves !A}$ निर्णेय नहीं है: हम विराम समस्या को इस
समस्या में घटाएँगे। जब तक हमारी चर्चा पूरी होगी, तब तक हम इससे कहीं अधिक
प्रबल परिणाम सिद्ध कर चुके होंगे।

$\Th{Q}$ की भाषा अंकगणित की भाषा है; $\Th{Q}$ में निम्न अभिगृहीत हैं
(जिन्हें !!{identity} सहित प्रथम-कोटि तर्कशास्त्र के अन्य अभिगृहीतों और
नियमों के साथ उपयोग करना है):
\begin{align*}
& \lforall[x][\lforall[y][(\eq[x'][y'] \lif \eq[x][y])]] \tag{$!Q_1$}\\
& \lforall[x][\eq/[\Obj 0][x']] \tag{$!Q_2$}\\
& \lforall[x][(\eq[x][\Obj 0] \lor \lexists[y][\eq[x][y']])] \tag{$!Q_3$}\\
& \lforall[x][\eq[(x + \Obj 0)][x]] \tag{$!Q_4$}\\
& \lforall[x][\lforall[y][\eq[(x + y')][(x + y)']]] \tag{$!Q_5$}\\
& \lforall[x][\eq[(x \times \Obj 0)][\Obj 0]] \tag{$!Q_6$}\\
& \lforall[x][\lforall[y][\eq[(x \times y')][((x \times y) + x)]]] \tag{$!Q_7$}\\
& \lforall[x][\lforall[y][(x < y \liff \lexists[z][\eq[(z' + x)][y]])]] \tag{$!Q_8$}
\end{align*}
प्रत्येक प्राकृतिक संख्या $n$ के लिए अंक-सूचक $\num{n}$ को वह पद
$\Obj{0}^{\prime\prime\ldots\prime}$ परिभाषित करें जिसमें कुल $n$ प्राइम
चिह्न हों। अतः $\num{0}$ केवल !!{constant}~$\Obj{0}$ है, $\num{1}$,
$\Obj{0}'$ है, $\num{2}$, $\Obj{0}''$ है, इत्यादि।

अंकगणित के सिद्धांत के रूप में $\Th{Q}$ \emph{अत्यंत} दुर्बल है; उदाहरणतः,
इसमें $\lforall[x][\eq/[x][x']]$ अथवा
$\lforall[x][\lforall[y][(x + y) = (y + x)]]$ जैसे बहुत सरल तथ्य भी सिद्ध
नहीं किए जा सकते। किंतु हम देखेंगे कि $\Th{Q}$ के इतना रोचक होने का बड़ा
कारण \emph{ठीक यही} है कि वह इतना दुर्बल है। वास्तव में वह अपूर्णता-प्रमेय को
लागू करने भर के लिए मुश्किल से पर्याप्त प्रबल है। $\Th{Q}$ के रोचक होने
का दूसरा कारण यह है कि उसके अभिगृहीतों का समुच्चय \emph{परिमित} है।

$\Th{Q}$ से अधिक प्रबल सिद्धांत (जिसे \emph{पियानो अंकगणित} $\Th{PA}$ कहते
हैं)~$\Th{Q}$ में आगमन की एक योजना जोड़कर मिलता है:
\[
(!A(\Obj 0) \land \lforall[x][(!A(x) \lif !A(x'))]) \lif \lforall[x][!A(x)]
\]
जहाँ $!A(x)$ कोई भी सूत्र है। यदि $!A(x)$ में $x$ के अतिरिक्त कोई और मुक्त
!!{variable}s हों, तो हम उन सबको बाँधने के लिए आरंभ में सार्विक परिमाणक
जोड़ते हैं (ताकि आगमन-योजना का संगत उदाहरण !!a{sentence} हो)। उदाहरणतः,
यदि $!A(x, y)$ में !!{variable}~$y$ भी मुक्त हो, तो संगत उदाहरण है
\[
\lforall[y][((!A(\Obj 0) \land \lforall[x][(!A(x) \lif !A(x'))]) \lif
  \lforall[x][!A(x)])]
\]
आगमन-योजना के उदाहरणों का उपयोग करके~$\Th{PA}$ के अभिगृहीतों से उन तथ्यों
से कहीं अधिक सिद्ध किया जा सकता है जो $\Th{Q}$ के अभिगृहीतों से सिद्ध होते
हैं। वास्तव में, प्राकृतिक संख्याओं के ऐसे ``स्वाभाविक'' कथन ढूँढ़ने में
काफी मेहनत लगती है जिन्हें पियानो अंकगणित में सिद्ध नहीं किया जा सकता!{}

\begin{defn}
\ollabel{defn:representable-fn}
  प्राकृतिक संख्याओं से प्राकृतिक संख्याओं में जाने वाला फलन
  $f(x_0,\ldots,x_k)$ {\em $\Th{Q}$ में निरूपणीय} कहलाता है यदि कोई सूत्र
  $!A_f(x_0,\dots,x_k,y)$ ऐसा हो कि जब भी $f(n_0,\dots,n_k) = m$ हो,
  $\Th{Q}$ निम्न दोनों को सिद्ध करे:
\begin{enumerate}
\item\ollabel{defn:rep:a} $!A_f(\num{n_0}, \dots, \num{n_k}, \num{m})$,
\item\ollabel{defn:rep:b} $\lforall[y][(!A_f(\num{n_0}, \dots,
\num{n_k}, y) \lif \num{m} = y)]$।
\end{enumerate}
\end{defn}

इस परिभाषा को कहने के और भी ढंग हैं; उदाहरणतः, हम समतुल्य रूप से यह माँग
सकते थे कि $\Th{Q}$,
$\lforall[y][(!A_f(\num{n_0}, \dots, \num{n_k}, y) \liff
\eq[y][\num{m}])]$ को सिद्ध करे।

\begin{thm}
\ollabel{thm:representable-iff-comp}
कोई फलन $\Th{Q}$ में तभी और केवल तभी निरूपणीय है जब वह संगणनीय है।
\end{thm}

इस प्रमेय के प्रमाण की दो दिशाएँ हैं। वाक्यविन्यास का अंकगणितीकरण उपलब्ध
हो जाने पर बाएँ-से-दाएँ दिशा काफी सीधी है। दूसरी दिशा के लिए अधिक काम
चाहिए। मूल विचार यह है: हम ``संगणनीय'' को सटीक बनाने के लिए ``सामान्य
पुनरावर्ती'' को चुनते हैं और दिखाते हैं कि प्रत्येक सामान्य पुनरावर्ती फलन
$\Th{Q}$ में निरूपणीय है। याद करें कि कोई फलन सामान्य पुनरावर्ती है यदि
उसे $\Zero$, उत्तराधिकारी फलन~$\Succ$ और प्रक्षेप फलनों~$\Proj{n}{i}$ से
संयोजन, आदिम पुनरावर्तन और नियमित न्यूनीकरण का उपयोग करके परिभाषित किया जा
सके। अतः प्रत्येक सामान्य पुनरावर्ती फलन को~$\Th{Q}$ में निरूपणीय दिखाने
का एक ढंग यह दिखाना है कि मूल फलन निरूपणीय हैं और जब कुछ फलन निरूपणीय हों,
तो उनसे संयोजन, आदिम पुनरावर्तन तथा नियमित न्यूनीकरण द्वारा परिभाषित फलन भी
निरूपणीय हों। दूसरे शब्दों में, हम दिखा सकते हैं कि मूल फलन निरूपणीय हैं
और निरूपणीय फलन संयोजन, आदिम पुनरावर्तन तथा नियमित न्यूनीकरण के अंतर्गत
``संवृत'' हैं। इससे प्रत्येक सामान्य पुनरावर्ती फलन की निरूपणीयता सुनिश्चित
होती है।

पर यह पता चलता है कि निरूपणीय फलनों का आदिम पुनरावर्तन के अंतर्गत संवृत
होना दिखाने वाला चरण कठिन है। इस चरण से बचने के लिए हम पहले दिखाते हैं कि
वास्तव में हम आदिम पुनरावर्तन के बिना काम चला सकते हैं। अर्थात् हम दिखाते
हैं कि प्रत्येक सामान्य पुनरावर्ती फलन को केवल संयोजन और नियमित न्यूनीकरण
का उपयोग करके मूल फलनों से परिभाषित किया जा सकता है। इसके लिए हम दिखाते
हैं कि आदिम पुनरावर्तन वस्तुतः एक विशिष्ट नियमित न्यूनीकरण द्वारा किया जा
सकता है। किंतु इसे सफल बनाने के लिए हमें कुछ अतिरिक्त मूल फलन जोड़ने होंगे:
जोड़, गुणन और सर्वसमिका-संबंध का अभिलाक्षणिक फलन~$\Char{=}$। तब हम यह
दिखाकर प्रमेय सिद्ध कर सकते हैं कि \emph{ये} सभी मूल फलन~$\Th{Q}$ में
निरूपणीय हैं और निरूपणीय फलन संयोजन तथा नियमित न्यूनीकरण के अंतर्गत संवृत
हैं।

\end{document}
